\section{Conclusion}\label{sec_conclusion}
LLMs trained via RL have exhibited emergent behaviors such as self-reflections, which offer a potential way for test-time scaling by backtracking to a previous state and generating a different action distribution with updated context. However, conventional RL training neither ensures the emergence of reflective exploration nor explains its benefits. The optimal RL policies are typically Markovian and have no incentive to enrich a revisited state with additional context since they depend on the history only through the state. Instead, RL exploration is confined to the training phase to identify action sequences that maximize cumulative rewards, and resorts to pure exploitation when the policy parameters are frozen after deployment. We propose to ground reflective exploration with a Bayesian RL objective that maximizes expected deployment-time return under a posterior distribution over MDPs induced by the training data. To maximize this objective, we propose BARL, which provides step-level guidance for when and how to engage in reflective exploration. BARL enables efficient exploration through hypothesis elimination and strategy switching. Our experiments are conducted on both synthetic and mathematical reasoning tasks, where we show that BARL outperforms conventional RL algorithms for evaluation benchmarks, and its exploration is more efficient. 
